% Modèle matOS : diaporama Beamer. Compiler avec pdflatex.
\documentclass[aspectratio=169,11pt]{beamer}
\usepackage[T1]{fontenc}
\usepackage[utf8]{inputenc}
\usepackage[french]{babel}
\usepackage{lmodern}
\usepackage{amsmath,amssymb,mathtools}

\usetheme{Madrid}
\usecolortheme{seahorse}
\setbeamertemplate{navigation symbols}{}
\newcommand{\R}{\mathbb{R}}

\title{Titre de l'exposé}
\subtitle{Séminaire de la mention Mathématiques}
\author{Prénom Nom}
\institute{Université d'Antananarivo}
\date{\today}

\begin{document}

\begin{frame}
  \titlepage
\end{frame}

\begin{frame}{Plan}
  \tableofcontents
\end{frame}

\section{Énoncé}

\begin{frame}{Théorème des accroissements finis}
  \begin{theorem}
    Soit $f : [a,b] \to \R$ continue sur $[a,b]$ et dérivable sur $]a,b[$.
    Il existe $c \in \,]a,b[$ tel que
    \[ f(b) - f(a) = f'(c)\,(b - a). \]
  \end{theorem}
  \pause
  \begin{block}{Idée de la preuve}
    Appliquer le théorème de Rolle à
    $g(x) = f(x) - \dfrac{f(b) - f(a)}{b - a}\,(x - a)$.
  \end{block}
\end{frame}

\section{Conséquences}

\begin{frame}{Conséquences}
  \begin{itemize}
    \item<1-> Une fonction de dérivée nulle sur un intervalle est constante.
    \item<2-> Si $|f'| \leq M$, alors $f$ est $M$-lipschitzienne.
  \end{itemize}
\end{frame}

\end{document}
